

\begin{frame}
    \frametitle{Espacios vectoriales}

    \[
        k^n,
        \qquad
        M_{m\times n}(k),
        \qquad
        k[x],
        \qquad
        k^X
    \]

    \begin{center}
        \alert{¿Qué tienen en común estos conjuntos?}
    \end{center}

    \medskip

    \begin{alertblock}{}
        En los primeros ejemplos puede pensarse en \(k=\mathbb{R}\).
    \end{alertblock}
\end{frame}


\begin{frame}
    \frametitle{El ejemplo prototípico: \(k^n\)}

    Dados
    \[
        u=(u_1,\ldots,u_n),
        \qquad
        v=(v_1,\ldots,v_n)\in k^n,
    \]
    y \(\lambda\in k\), definimos
    \[
        u+v
        =
        (u_1+v_1,\ldots,u_n+v_n)
    \]
    y
    \[
        \lambda u
        =
        (\lambda u_1,\ldots,\lambda u_n).
    \]

    Estas operaciones satisfacen, entre otras, las propiedades
    \[
        u+v=v+u,
        \qquad
        (u+v)+w=u+(v+w),
    \]
    \[
        \lambda(u+v)=\lambda u+\lambda v,
        \qquad
        (\lambda+\mu)u=\lambda u+\mu u.
    \]
\end{frame}


\begin{frame}
    \frametitle{Otros ejemplos}

    Las mismas operaciones pueden definirse en conjuntos de naturaleza muy distinta:

    \begin{itemize}
        \item las matrices
        \[
            M_{m\times n}(k),
        \]

        \item los polinomios
        \[
            k[x],
        \]

        \item las funciones de un conjunto \(X\) en \(k\)
        \[
            k^X=\{f\colon X\longrightarrow k\}.
        \]
    \end{itemize}

    En todos estos casos tenemos una suma y un producto por escalares
    que satisfacen las mismas propiedades formales que en \(k^n\).

    \medskip

    \begin{center}
        \alert{La naturaleza de los elementos cambia, pero la estructura es la misma.}
    \end{center}
\end{frame}

\begin{frame}
    \frametitle{Espacios vectoriales}

    \begin{definition}
        Sea \(k\) un cuerpo. Un \alert{espacio vectorial sobre \(k\)} es un conjunto \(V\)
        dotado de dos operaciones
        \begin{alignat*}{2}
            +\colon V\times V &\longrightarrow V, \qquad &(u,v) &\longmapsto u+v,\\
            \cdot\;\colon k\times V &\longrightarrow V, \qquad &(\lambda,v) &\longmapsto \lambda v,
        \end{alignat*}
        llamadas, respectivamente, \alert{suma} y \alert{producto por escalares},
        que satisfacen ciertos axiomas.
    \end{definition}

    \medskip

    Los elementos de \(V\) se llaman \alert{vectores}.
\end{frame}


\begin{frame}
    \frametitle{Axiomas de un espacio vectorial}

    Para cualesquiera \(u,v,w\in V\),

    \begin{align*}
        (u+v)+w &= u+(v+w),\\
        u+v &= v+u,
    \end{align*}

    existe un elemento \(0\in V\) llamado \alert{vector cero}, \alert{nulo} o \alert{trivial} tal que
    \[
        v+0=v,
    \]

    y para cada \(v\in V\) existe un elemento \(-v\in V\) tal que
    \[
        v+(-v)=0.
    \]

    \medskip

    Es decir, \((V,+)\) es un \alert{grupo abeliano}.
\end{frame}


\begin{frame}
    \frametitle{Axiomas de un espacio vectorial}

    Para cualesquiera \(\lambda,\mu\in k\) y \(u,v\in V\),

    \begin{align*}
        \lambda(u+v) &= \lambda u+\lambda v,\\
        (\lambda+\mu)v &= \lambda v+\mu v,\\
        (\lambda\mu)v &= \lambda(\mu v),\\
        1v &= v.
    \end{align*}

    \medskip

    Obsérvese que utilizamos el mismo símbolo \(0\) para denotar
    el elemento neutro de \(k\) y el vector cero de \(V\):
    \[
        0\in k,
        \qquad
        0\in V.
    \]
\end{frame}


\begin{frame}
    \frametitle{El espacio vectorial \(k^n\)}

    \begin{example}
        El conjunto
        \[
            k^n=\{(u_1,\ldots,u_n):u_i\in k\}
        \]
        es un espacio vectorial sobre \(k\) con las operaciones
        \[
            (u_1,\ldots,u_n)+(v_1,\ldots,v_n)
            =
            (u_1+v_1,\ldots,u_n+v_n),
        \]
        \[
            \lambda(u_1,\ldots,u_n)
            =
            (\lambda u_1,\ldots,\lambda u_n).
        \]
    \end{example}

    \medskip

    Su vector cero es
    \[
        (0,\ldots,0),
    \]
    y el opuesto de \((u_1,\ldots,u_n)\) es
    \[
        (-u_1,\ldots,-u_n).
    \]
\end{frame}


\begin{frame}
    \frametitle{El espacio vectorial de las matrices}

    \begin{example}
        El conjunto
        \[
            M_{m\times n}(k)
        \]
        de matrices \(m\times n\) con coeficientes en \(k\) es un espacio vectorial
        sobre \(k\), con la suma de matrices y el producto por escalares habituales.
    \end{example}

    \medskip

    El vector cero es la matriz nula
    \[
        0=
        \begin{pmatrix}
            0 & \cdots & 0\\
            \vdots & & \vdots\\
            0 & \cdots & 0
        \end{pmatrix},
    \]
    y el opuesto de una matriz \(A\) es \(-A\).
\end{frame}


\begin{frame}
    \frametitle{El espacio vectorial de los polinomios}

    \begin{example}
        El conjunto
        \[
            k[x]
        \]
        de polinomios con coeficientes en \(k\) es un espacio vectorial sobre \(k\),
        con la suma y el producto por escalares habituales.
    \end{example}

    \medskip

    También lo es el conjunto
    \[
        k_n[x]
        =
        \{p(x)\in k[x]:\deg p\leq n\}\cup\{0\}.
    \]

    \medskip

    Por ejemplo,
    \[
        1,\quad x,\quad x^2,\quad \ldots,\quad x^n
    \]
    son vectores de \(k_n[x]\).
\end{frame}


\begin{frame}
    \frametitle{Espacios vectoriales de funciones}

    \begin{example}
        Sea \(X\) un conjunto. El conjunto
        \[
            k^X=\{f\colon X\longrightarrow k\}
        \]
        es un espacio vectorial sobre \(k\), definiendo
        \[
            (f+g)(x)=f(x)+g(x),
        \]
        \[
            (\lambda f)(x)=\lambda f(x),
        \]
        para todo \(x\in X\).
    \end{example}

    \medskip

    El vector cero es la función constante
    \[
        0\colon X\longrightarrow k,
        \qquad
        x\longmapsto 0.
    \]
\end{frame}

\begin{frame}
    \frametitle{Un ejemplo menos evidente}

    \begin{example}
        Consideremos
        \[
            V=\mathbb{R}_{>0}
        \]
        y definamos la \alert{suma} y el \alert{producto por escalares} del siguiente modo inusual:
        \[
            x\oplus y=xy,
            \qquad
            \lambda\odot x=x^\lambda,
        \]
        para \(x,y\in V\) y \(\lambda\in\mathbb{R}\). Entonces \(V\) es un espacio vectorial sobre \(\mathbb{R}\). En este espacio vectorial, el vector cero es
    \[
        1,
    \]
    pues
    \[
        x\oplus 1=x.
    \]

    Además, el opuesto de \(x\) es
    \[
        x^{-1}.
    \]
        \end{example}
\end{frame}

\begin{frame}
    \frametitle{Resta de vectores}

    Sean \(u,v\in V\). Definimos la \alert{resta} de \(v\) a \(u\) por
    \[
        u-v:=u+(-v).
    \]

    \medskip

    Es decir, restar un vector es, por definición, sumar su opuesto.

    \medskip

    En particular,
    \[
        u-v=w
        \quad\Longleftrightarrow\quad
        u=w+v,
    \]
    es decir, podemos despejar.
\end{frame}

\begin{frame}
    \frametitle{Unicidad del vector cero}

    \begin{proposition}
        Sea \(V\) un espacio vectorial sobre \(k\). El vector cero de \(V\) es único.
    \end{proposition}

    \begin{proof}
        Supongamos que \(0,0'\in V\) son dos vectores cero. Entonces
        \[
            0=0+0'=0'.
        \]
        Por tanto, el vector cero es único.
    \end{proof}
\end{frame}

\begin{frame}
    \frametitle{Unicidad del opuesto}

    \begin{proposition}
        Sea \(V\) un espacio vectorial sobre \(k\). Para cada \(v\in V\),
        existe un único vector \(-v\in V\) tal que
        \[
            v+(-v)=0.
        \]
    \end{proposition}

    \begin{proof}
        La existencia forma parte de los axiomas de espacio vectorial.

        Supongamos que \(w,w'\in V\) satisfacen
        \[
            v+w=0,
            \qquad
            v+w'=0.
        \]
        Entonces
        \[
            w
            =
            w+0
            =
            w+(v+w')
            =
            (w+v)+w'
            =
            (v+w)+w'
            =
            0+w'
            =
            w'+0
            =
            w'.
        \]

        Por tanto, el opuesto de \(v\) es único.
    \end{proof}
\end{frame}

\begin{frame}[allowframebreaks, t]
    \frametitle{Primeras consecuencias de los axiomas}

    \begin{proposition}
        Sea \(V\) un espacio vectorial sobre \(k\). Para todo \(v\in V\)
        y todo \(\lambda\in k\),
        \[
            0v=0,
            \qquad
            \lambda 0=0,
            \qquad
            (-1)v=-v.
        \]
    \end{proposition}

    \begin{proof}
        Como
        \[
            0v=(0+0)v=0v+0v,
        \]
        despejando obtenemos
        \[
            0v=0.
        \]

        Análogamente,
        \[
            \lambda 0
            =
            \lambda(0+0)
            =
            \lambda 0+\lambda 0,
        \]
        y por tanto
        \[
            \lambda 0=0.
        \]

        Por otra parte,
        \[
            v+(-1)v
            =
            1v+(-1)v
            =
            (1-1)v
            =
            0v
            =
            0.
        \]
        Por unicidad del opuesto,
        \[
            (-1)v=-v.
        \]
    \end{proof}
\end{frame}

\begin{frame}
    \frametitle{Despejar un escalar no nulo}

    Supongamos que
    \[
        \lambda u=v,
    \]
    con \(\lambda\in k\) no nulo. Como \(k\) es un cuerpo, podemos dividir por \(\lambda\), es decir, existe \(\lambda^{-1}\in k\) tal que
    \[
        \lambda^{-1}\lambda=1.
    \]
    Multiplicando los dos miembros de la primera ecuación por \(\lambda^{-1}\),
    \[
        \lambda^{-1}(\lambda u)=\lambda^{-1}v.
    \]
    Por los axiomas del producto por escalares,
    \[
        \lambda^{-1}(\lambda u)=(\lambda^{-1}\lambda)u=1u=u,
    \]
    y por tanto
    \[
        u=\lambda^{-1}v.
    \]
    Es decir, podemos despejar escalares no nulos de una ecuación vectorial.
\end{frame}

\begin{frame}
    \frametitle{Cuando el producto por escalares es cero}

    \begin{corollary}
        Sea \(V\) un espacio vectorial sobre \(k\). Si \(\lambda\in k\) y \(v\in V\),
        entonces
        \[
            \lambda v=0
            \quad\Longrightarrow\quad
            \lambda=0
            \quad\text{o}\quad
            v=0.
        \]
    \end{corollary}

    \begin{proof}
        Si \(\lambda\neq 0\), entonces 
        \[v=\lambda^{-1}0=0.\]
    \end{proof}
\end{frame}


\begin{frame}
    \frametitle{Un conjunto que no es un espacio vectorial}

    Consideremos
    \[
        V=\{(x,y)\in\mathbb{R}^2:x\geq 0\}
    \]
    con las operaciones habituales de \(\mathbb{R}^2\).

    \medskip

    ¿Es \(V\) un espacio vectorial sobre \(\mathbb{R}\)?

    \medskip

    No. Por ejemplo,
    \[
        (1,0)\in V,
    \]
    pero
    \[
        -(1,0)=(-1,0)\notin V.
    \]

    \medskip

    Por tanto, \(V\) no es cerrado respecto de la toma de opuestos.
\end{frame}


\begin{frame}
    \frametitle{Otro conjunto que no es un espacio vectorial}

    Consideremos el conjunto de polinomios de grado exactamente \(n\):
    \[
        V=\{p(x)\in k[x]:\deg p=n\}.
    \]

    \medskip

    ¿Es \(V\) un espacio vectorial sobre \(k\)?

    \medskip

    No. El polinomio cero no pertenece a \(V\).

\end{frame}


\begin{frame}
    \frametitle{Operaciones que no definen un espacio vectorial}

    En \(\mathbb{R}^2\), consideremos la suma habitual y definamos
    \[
        \lambda\cdot(x,y)=(\lambda x,y).
    \]

    \medskip

    ¿Definen estas operaciones una estructura de espacio vectorial sobre \(\mathbb{R}\)?

    \medskip

    No. Por ejemplo,
    \[
        0\cdot(x,y)=(0,y),
    \]
    que en general no es el vector cero.

    \medskip

    Por tanto, falla una consecuencia necesaria de los axiomas:
    \[
        0v=0.
    \]
\end{frame}


\begin{frame}
    \frametitle{Ejercicio}

    Determina cuáles de los siguientes conjuntos, con las operaciones habituales,
    son espacios vectoriales sobre \(k\):

    \begin{enumerate}
        \item
        \(
            \{(x_1,\ldots,x_n)\in k^n:
            x_1+\cdots+x_n=0\}.
        \)

        \item
        \(
            \{p(x)\in k[x]:\deg p=n\}.
        \)

        \item
        \(
            \{p(x)\in k[x]:p(0)=0\}.
        \)
    \end{enumerate}

    \medskip

    En cada caso, justifica la respuesta a partir de los axiomas.
\end{frame}